
% https://akaedu.github.io/book/ch18.html

\chapter{x86汇编程序基础}
\BookAnchor{db.ch18}\BookAnchor{db.ch18.id2769122}%
要彻底搞清楚C语言的原理，就必须深入到指令一层去理解。你写一行C代码，编译器会生成什么样的指令，要做到心中有数。本章介绍汇编程序的一些基础知识。汇编不是本书的重点，本书要求读者能看懂基本的汇编程序而不要求会写汇编程序，下一章将在汇编的基础上讨论C语言的原理。

\BookNav{\hyperlink{db.ch17s05}{上一页}\enspace\textperiodcentered\enspace \hyperlink{db.pt02}{上一级}\enspace\textperiodcentered\enspace \hyperlink{db.index}{起始页}\enspace\textperiodcentered\enspace \hyperlink{db.ch18s01}{下一页}}


% https://akaedu.github.io/book/ch18s01.html

\section{最简单的汇编程序}
\BookAnchor{db.ch18s01}\BookAnchor{db.ch18s01.id2769131}%
\BookExampleTitle{最简单的汇编程序}\BookAnchor{db.ch18s01.id2769136}

\BookSource{numbers=left,xleftmargin=19pt}{code/ch18s01/01.pv}{9}{33}


把这个程序保存成文件\BookCode{hello.\BookSplit{}s}（汇编程序通常以\BookCode{.\BookSplit{}s}作为文件名后缀），用汇编器（Assembler）\BookAnchor{db.ch18s01.id2769149}\BookCode{as}把汇编程序中的助记符翻译成机器指令，生成目标文件\BookCode{hello.\BookSplit{}o}：


\BookTerminal{code/ch18s01/02.txt}{3}{1}


然后用链接器（Linker，或Link Editor）\BookAnchor{db.ch18s01.id2769187}\BookCode{ld}把目标文件\BookCode{hello.\BookSplit{}o}链接成可执行文件\BookCode{hello}：


\BookTerminal{code/ch18s01/03.txt}{3}{1}


为什么用汇编器翻译成机器指令了还不行，还要有一个链接的步骤呢？链接主要有两个作用，一是修改目标文件中的信息，对地址做重定位，在\hyperlink{db.ch18s05.id2771639}{第 5.\BookSplit{}2 节 “可执行文件”}详细解释，二是把多个目标文件合并成一个可执行文件，在\hyperlink{db.ch19s02.id2776173}{第 2 节 “\BookCode{main}函数和启动例程”}详细解释。我们这个例子虽然只有一个目标文件，但也需要经过链接才能成为可执行文件。



现在执行这个程序，它只做了一件事就是退出，退出状态是4，\hyperlink{db.ch03s02.id2712975}{第 2 节 “自定义函数”}讲过在Shell中可以用特殊变量\BookCode{\$?}得到上一条命令的退出状态：


\BookTerminal{code/ch18s01/04.txt}{5}{3}


所以这段汇编代码相当于在C程序的\BookCode{main}函数中\BookCode{return 4;\BookSplit{}}。为什么会相当呢？我们在\hyperlink{db.ch19s02.id2776173}{第 2 节 “\BookCode{main}函数和启动例程”}详细解释。



下面逐行分析这个汇编程序。首先，\BookCode{\#}号表示单行注释，类似于C语言的\BookCode{/\BookSplit{}/\BookSplit{}}注释。


\BookSource{numbers=none}{code/ch18s01/05.pv}{3}{1}


汇编程序中以\BookCode{.\BookSplit{}}开头的名称并不是指令的助记符，不会被翻译成机器指令，而是给汇编器一些特殊指示，称为汇编指示（Assembler Directive）\BookAnchor{db.ch18s01.id2769314}或伪操作（Pseudo-operation）\BookAnchor{db.ch18s01.id2769324}，由于它不是真正的指令所以加个“伪”字。\BookCode{.\BookSplit{}section}指示把代码划分成若干个段（Section）\BookAnchor{db.ch18s01.id2769345}，程序被操作系统加载执行时，每个段被加载到不同的地址，操作系统对不同的页面设置不同的读、写、执行权限。\BookCode{.\BookSplit{}data}段保存程序的数据，是可读可写的，相当于C程序的全局变量。本程序中没有定义数据，所以\BookCode{.\BookSplit{}data}段是空的。


\BookSource{numbers=none}{code/ch18s01/06.pv}{3}{1}


\BookCode{.\BookSplit{}text}段保存代码，是只读和可执行的，后面那些指令都属于\BookCode{.\BookSplit{}text}段。


\BookSource{numbers=none}{code/ch18s01/07.pv}{3}{1}


\BookCode{\_\BookSplit{}start}是一个符号（Symbol）\BookAnchor{db.ch18s01.id2769400}，符号在汇编程序中代表一个地址，可以用在指令中，汇编程序经过汇编器的处理之后，所有的符号都被替换成它所代表的地址值。在C语言中我们通过变量名访问一个变量，其实就是读写某个地址的内存单元，我们通过函数名调用一个函数，其实就是跳转到该函数第一条指令所在的地址，所以变量名和函数名都是符号，本质上是代表内存地址的。



\BookCode{.\BookSplit{}globl}指示告诉汇编器，\BookCode{\_\BookSplit{}start}这个符号要被链接器用到，所以要在目标文件的符号表中标记它是一个全局符号（在\hyperlink{db.ch18s05.id2770854}{第 5.\BookSplit{}1 节 “目标文件”}详细解释）。\BookCode{\_\BookSplit{}start}就像C程序的\BookCode{main}函数一样特殊，是整个程序的入口，链接器在链接时会查找目标文件中的\BookCode{\_\BookSplit{}start}符号代表的地址，把它设置为整个程序的入口地址，所以每个汇编程序都要提供一个\BookCode{\_\BookSplit{}start}符号并且用\BookCode{.\BookSplit{}globl}声明。如果一个符号没有用\BookCode{.\BookSplit{}globl}声明，就表示这个符号不会被链接器用到。


\BookPlain{code/ch18s01/08.txt}{3}{1}


这里定义了\BookCode{\_\BookSplit{}start}符号，汇编器在翻译汇编程序时会计算每个数据对象和每条指令的地址，当看到这样一个符号定义时，就把它后面一条指令的地址作为这个符号所代表的地址。而\BookCode{\_\BookSplit{}start}这个符号又比较特殊，它所代表的地址是整个程序的入口地址，所以下一条指令\BookCode{movl \$1,\BookSplit{} \%eax}就成了程序中第一条被执行的指令。


\BookSource{numbers=none}{code/ch18s01/09.pv}{3}{1}


这是一条数据传送指令，这条指令要求CPU内部产生一个数字1并保存到\BookCode{eax}寄存器中。\BookCode{mov}的后缀l表示long，说明是32位的传送指令。这条指令不要求CPU读内存，1这个数是在CPU内部产生的，称为立即数（Immediate）\BookAnchor{db.ch18s01.id2769538}。在汇编程序中，立即数前面要加\$，寄存器名前面要加\%，以便跟符号名区分开。以后我们会看到\BookCode{mov}指令还有另外几种形式，但数据传送方向都是一样的，第一个操作数总是源操作数，第二个操作数总是目标操作数。


\BookSource{numbers=none}{code/ch18s01/10.pv}{3}{1}


和上一条指令类似，生成一个立即数4并保存到\BookCode{ebx}寄存器中。


\BookPlain{code/ch18s01/11.txt}{3}{1}


前两条指令都是为这条指令做准备的，执行这条指令时发生以下动作：


\begin{enumerate}

\item \BookCode{int}指令称为软中断指令，可以用这条指令故意产生一个异常，上一章讲过，异常的处理和中断类似，CPU从用户模式切换到特权模式，然后跳转到内核代码中执行异常处理程序。

\item \BookCode{int}指令中的立即数0x80是一个参数，在异常处理程序中要根据这个参数决定如何处理，在Linux内核中\BookCode{int \$0x80}这种异常称为系统调用（System Call）\BookAnchor{db.ch18s01.id2769615}。内核提供了很多系统服务供用户程序使用，但这些系统服务不能像库函数（比如\BookCode{printf}）那样调用，因为在执行用户程序时CPU处于用户模式，不能直接调用内核函数，所以需要通过系统调用切换CPU模式，经由异常处理程序进入内核，用户程序只能通过寄存器传几个参数，之后就要按内核设计好的代码路线走，而不能由用户程序随心所欲，想调哪个内核函数就调哪个内核函数，这样可以保证系统服务被安全地调用。在调用结束之后，CPU再切换回用户模式，继续执行\BookCode{int \$0x80}的下一条指令，在用户程序看来就像函数调用和返回一样。

\item \BookCode{eax}和\BookCode{ebx}的值是传递给系统调用的两个参数。\BookCode{eax}的值是系统调用号，Linux的各种系统调用都是由\BookCode{int \$0x80}指令引发的，内核需要通过\BookCode{eax}判断用户要调哪个系统调用，\BookCode{\_\BookSplit{}exit}的系统调用号是1。\BookCode{ebx}的值是传给\BookCode{\_\BookSplit{}exit}的参数，表示退出状态。大多数系统调用完成之后会返回用户空间继续执行后面的指令，而\BookCode{\_\BookSplit{}exit}系统调用比较特殊，它会终止掉当前进程，而不是返回用户空间继续执行。
\end{enumerate}

\begin{BookNote}
\BookNoteTitle{x86汇编的两种语法：intel语法和AT\&T语法}


x86汇编一直存在两种不同的语法，在intel的官方文档中使用intel语法，Windows也使用intel语法，而UNIX平台的汇编器一直使用AT\&T语法，所以本书使用AT\&T语法。\BookCode{movl \%edx,\BookSplit{}\%eax}这条指令如果用intel语法来写，就是\BookCode{mov eax,\BookSplit{}edx}，寄存器名不加\%号，源操作数和目标操作数的位置互换，字长也不是用指令的后缀l表示而是用另外的方式表示。本书不详细讨论这两种语法之间的区别，读者可以参考\hyperlink{db.bi01.bibli.assemblyhowto}{[AssemblyHOWTO]}。



介绍x86汇编的书很多，UNIX平台的书都采用AT\&T语法，例如\hyperlink{db.bi01.bibli.groundup}{[GroudUp]}，其它书一般采用intel语法，例如\hyperlink{db.bi01.bibli.x86assembly}{[x86Assembly]}。

\end{BookNote}

\subsection*{习题}
\BookAnchor{db.ch18s01.id2769776}%


1、把本节例子中的\BookCode{int \$0x80}指令去掉，汇编、链接也能通过，但是执行的时候出现段错误，你能解释其原因吗？

\BookNav{\hyperlink{db.ch18}{上一页}\enspace\textperiodcentered\enspace \hyperlink{db.ch18}{上一级}\enspace\textperiodcentered\enspace \hyperlink{db.index}{起始页}\enspace\textperiodcentered\enspace \hyperlink{db.ch18s02}{下一页}}


% https://akaedu.github.io/book/ch18s02.html

\section{x86的寄存器}
\BookAnchor{db.ch18s02}\BookAnchor{db.ch18s02.id2769794}%
x86的通用寄存器有\BookCode{eax}、\BookCode{ebx}、\BookCode{ecx}、\BookCode{edx}、\BookCode{edi}、\BookCode{esi}。这些寄存器在大多数指令中是可以任意选用的，比如\BookCode{movl}指令可以把一个立即数传送到\BookCode{eax}中，也可传送到\BookCode{ebx}中。但也有一些指令规定只能用其中某个寄存器做某种用途，例如除法指令\BookCode{idivl}要求被除数在\BookCode{eax}寄存器中，\BookCode{edx}寄存器必须是0，而除数可以在任意寄存器中，计算结果的商数保存在\BookCode{eax}寄存器中（覆盖原来的被除数），余数保存在\BookCode{edx}寄存器中。也就是说，通用寄存器对于某些特殊指令来说也不是通用的。



x86的特殊寄存器有\BookCode{ebp}、\BookCode{esp}、\BookCode{eip}、\BookCode{eflags}。\BookCode{eip}是程序计数器，\BookCode{eflags}保存着计算过程中产生的标志位，其中包括\hyperlink{db.ch14s03.id2753602}{第 3 节 “整数的加减运算”}讲过的进位标志、溢出标志、零标志和负数标志，在intel的手册中这几个标志位分别称为CF、OF、ZF、SF。\BookCode{ebp}和\BookCode{esp}用于维护函数调用的栈帧，在\hyperlink{db.ch19s01.id2774377}{第 1 节 “函数调用”}详细讨论。

\BookNav{\hyperlink{db.ch18s01}{上一页}\enspace\textperiodcentered\enspace \hyperlink{db.ch18}{上一级}\enspace\textperiodcentered\enspace \hyperlink{db.index}{起始页}\enspace\textperiodcentered\enspace \hyperlink{db.ch18s03}{下一页}}


% https://akaedu.github.io/book/ch18s03.html

\section{第二个汇编程序}
\BookAnchor{db.ch18s03}\BookAnchor{db.ch18s03.id2769960}%
\BookExampleTitle{求一组数的最大值的汇编程序}\BookAnchor{db.ch18s03.id2769965}

\BookSource{numbers=left,xleftmargin=19pt}{code/ch18s03/01.pv}{9}{42}


汇编、链接、运行：


\BookTerminal{code/ch18s03/02.txt}{7}{5}


这个程序在一组数中找到一个最大的数，并把它作为程序的退出状态。这组数在\BookCode{.\BookSplit{}data}段给出：


\BookPlain{code/ch18s03/03.txt}{4}{2}


\BookCode{.\BookSplit{}long}指示声明一组数，每个数占32位，相当于C语言中的数组。这个数组开头定义了一个符号\BookCode{data\_\BookSplit{}items}，汇编器会把数组的首地址作为\BookCode{data\_\BookSplit{}items}符号所代表的地址，\BookCode{data\_\BookSplit{}items}类似于C语言中的数组名。\BookCode{data\_\BookSplit{}items}这个标号没有用\BookCode{.\BookSplit{}globl}声明，因为它只在这个汇编程序内部使用，链接器不需要用到这个名字。除了\BookCode{.\BookSplit{}long}之外，常用的数据声明还有：


\begin{itemize}

\item \BookCode{.\BookSplit{}byte}，也是声明一组数，每个数占8位

\item \BookCode{.\BookSplit{}ascii}，例如\BookCode{.\BookSplit{}ascii \BookQuoteChar{}Hello w\BookSplit{}orld\BookQuoteChar{}}，声明11个数，取值为相应字符的ASCII码。注意，和C语言不同，这样声明的字符串末尾是没有\BookCode{'\textbackslash{}0'}字符的，如果需要以\BookCode{'\textbackslash{}0'}结尾可以声明为\BookCode{.\BookSplit{}ascii \BookQuoteChar{}Hello w\BookSplit{}orld\textbackslash{}0\BookQuoteChar{}}。
\end{itemize}


\BookCode{data\_\BookSplit{}items}数组的最后一个数是0，我们在一个循环中依次比较每个数，碰到0的时候让循环终止。在这个循环中：


\begin{itemize}

\item \BookCode{edi}寄存器保存数组中的当前位置，每次比较完一个数就把\BookCode{edi}的值加1，指向数组中的下一个数。

\item \BookCode{ebx}寄存器保存到目前为止找到的最大值，如果发现有更大的数就更新\BookCode{ebx}的值。

\item \BookCode{eax}寄存器保存当前要比较的数，每次更新\BookCode{edi}之后，就把下一个数读到\BookCode{eax}中。
\end{itemize}

\BookSource{numbers=none}{code/ch18s03/04.pv}{4}{2}


初始化\BookCode{edi}，指向数组的第0个元素。


\BookSource{numbers=none}{code/ch18s03/05.pv}{3}{1}


这条指令把数组的第0个元素传送到\BookCode{eax}寄存器中。\BookCode{data\_\BookSplit{}items}是数组的首地址，\BookCode{edi}的值是数组的下标，4表示数组的每个元素占4字节，那么数组中第\BookCode{edi}个元素的地址应该是\BookCode{data\_\BookSplit{}items +\BookSplit{} edi * 4}，写在指令中就是\BookCode{data\_\BookSplit{}items(\BookSplit{},\BookSplit{}\%edi,\BookSplit{}4)\BookSplit{}}，这种地址表示方式在下一节还会详细解释。


\BookSource{numbers=none}{code/ch18s03/06.pv}{3}{1}


\BookCode{ebx}的初始值也是数组的第0个元素。下面我们进入一个循环，循环的开头定义一个符号\BookCode{start\_\BookSplit{}loop}，循环的末尾之后定义一个符号\BookCode{loop\_\BookSplit{}exit}。


\BookSource{numbers=none}{code/ch18s03/07.pv}{5}{3}


比较\BookCode{eax}的值是不是0，如果是0就说明到达数组末尾了，就要跳出循环。\BookCode{cmpl}指令将两个操作数相减，但计算结果并不保存，只是根据计算结果改变\BookCode{eflags}寄存器中的标志位。如果两个操作数相等，则计算结果为0，\BookCode{eflags}中的ZF位置1。\BookCode{je}是一个条件跳转指令，它检查\BookCode{eflags}中的ZF位，ZF位为1则发生跳转，ZF位为0则不跳转，继续执行下一条指令。可见比较指令和条件跳转指令是配合使用的，前者改变标志位，后者根据标志位决定是否跳转。\BookCode{je}可以理解成“jump if equal”，如果参与比较的两数相等则跳转。


\BookSource{numbers=none}{code/ch18s03/08.pv}{4}{2}


将\BookCode{edi}的值加1，把数组中的下一个数传送到\BookCode{eax}寄存器中。


\BookSource{numbers=none}{code/ch18s03/09.pv}{4}{2}


把当前数组元素\BookCode{eax}和目前为止找到的最大值\BookCode{ebx}做比较，如果前者小于等于后者，则最大值没有变，跳转到循环开头比较下一个数，否则继续执行下一条指令。\BookCode{jle}表示“jump if less than or equal”。


\BookSource{numbers=none}{code/ch18s03/10.pv}{4}{2}


更新了最大值\BookCode{ebx}然后跳转到循环开头比较下一个数。\BookCode{jmp}是一个无条件跳转指令，什么条件也不判断，直接跳转。\BookCode{loop\_\BookSplit{}exit}符号后面的指令调\BookCode{\_\BookSplit{}exit}系统调用退出程序。

\BookNav{\hyperlink{db.ch18s02}{上一页}\enspace\textperiodcentered\enspace \hyperlink{db.ch18}{上一级}\enspace\textperiodcentered\enspace \hyperlink{db.index}{起始页}\enspace\textperiodcentered\enspace \hyperlink{db.ch18s04}{下一页}}


% https://akaedu.github.io/book/ch18s04.html

\section{寻址方式}
\BookAnchor{db.ch18s04}\BookAnchor{db.ch18s04.id2770425}%
通过上一节的例子我们了解到，访问内存时在指令中可以用多种方式表示内存地址，比如可以用数组基地址、元素长度和下标三个量来表示，增加了寻址的灵活性。本节介绍x86常用的几种寻址方式（Addressing Mode）\BookAnchor{db.ch18s04.id2770436}。内存寻址在指令中可以表示成如下的通用格式：


\begin{BookLiteral}
\BookLiteralLine{ADDRESS\_OR\_OFFSET(\%BASE\_OR\_OFFSET,\%INDEX,MULTIPLIER)}
\end{BookLiteral}


它所表示的地址可以这样计算出来：


\begin{BookLiteral}
\BookLiteralLine{FINAL ADDRESS = ADDRESS\_OR\_OFFSET + BASE\_OR\_OFFSET + MULTIPLIER * INDEX}
\end{BookLiteral}


其中ADDRESS\_\BookSplit{}OR\_\BookSplit{}OFFSET和MULTIPLIER必须是常数，BASE\_\BookSplit{}OR\_\BookSplit{}OFFSET和INDEX必须是寄存器。在有些寻址方式中会省略这4项中的某些项，相当于这些项是0。


\begin{itemize}

\item 直接寻址（Direct Addressing Mode）\BookAnchor{db.ch18s04.id2770476}。只使用ADDRESS\_\BookSplit{}OR\_\BookSplit{}OFFSET寻址，例如\BookCode{movl ADDRESS,\BookSplit{} \%eax}把ADDRESS地址处的32位数传送到\BookCode{eax}寄存器。

\item 变址寻址（Indexed Addressing Mode）\BookAnchor{db.ch18s04.id2770501} 。上一节的\BookCode{movl data\_\BookSplit{}items(\BookSplit{},\BookSplit{}\%edi,\BookSplit{}4)\BookSplit{},\BookSplit{} \%eax}就属于这种寻址方式，用于访问数组元素比较方便。

\item 间接寻址（Indirect Addressing Mode）\BookAnchor{db.ch18s04.id2770521}。只使用BASE\_\BookSplit{}OR\_\BookSplit{}OFFSET寻址，例如\BookCode{movl (\BookSplit{}\%eax)\BookSplit{},\BookSplit{} \%ebx}，把\BookCode{eax}寄存器的值看作地址，把内存中这个地址处的32位数传送到\BookCode{ebx}寄存器。注意和\BookCode{movl \%eax,\BookSplit{} \%ebx}区分开。

\item 基址寻址（Base Pointer Addressing Mode）\BookAnchor{db.ch18s04.id2770560}。只使用ADDRESS\_\BookSplit{}OR\_\BookSplit{}OFFSET和BASE\_\BookSplit{}OR\_\BookSplit{}OFFSET寻址，例如\BookCode{movl 4(\BookSplit{}\%eax)\BookSplit{},\BookSplit{} \%ebx}，用于访问结构体成员比较方便，例如一个结构体的基地址保存在\BookCode{eax}寄存器中，其中一个成员在结构体内的偏移量是4字节，要把这个成员读上来就可以用这条指令。

\item 立即数寻址（Immediate Mode）\BookAnchor{db.ch18s04.id2770588}。就是指令中有一个操作数是立即数，例如\BookCode{movl \$12,\BookSplit{} \%eax}中的\BookCode{\$12}，这其实跟寻址没什么关系，但也算作一种寻址方式。

\item 寄存器寻址（Register Addressing Mode）\BookAnchor{db.ch18s04.id2770614}。就是指令中有一个操作数是寄存器，例如\BookCode{movl \$12,\BookSplit{} \%eax}中的\BookCode{\%eax}，这跟内存寻址没什么关系，但也算作一种寻址方式。在汇编程序中寄存器用助记符来表示，在机器指令中则要用几个Bit表示寄存器的编号，这几个Bit也可以看作寄存器的地址，但是和内存地址不在一个地址空间。
\end{itemize}

\BookNav{\hyperlink{db.ch18s03}{上一页}\enspace\textperiodcentered\enspace \hyperlink{db.ch18}{上一级}\enspace\textperiodcentered\enspace \hyperlink{db.index}{起始页}\enspace\textperiodcentered\enspace \hyperlink{db.ch18s05}{下一页}}


% https://akaedu.github.io/book/ch18s05.html

\section{ELF文件}
\BookAnchor{db.ch18s05}\BookAnchor{db.ch18s05.id2770642}%
ELF文件格式是一个开放标准，各种UNIX系统的可执行文件都采用ELF格式，它有三种不同的类型：


\begin{itemize}

\item 可重定位的目标文件（Relocatable，或者Object File）\BookAnchor{db.ch18s05.id2770661}\BookAnchor{db.ch18s05.id2770669}

\item 可执行文件（Executable）\BookAnchor{db.ch18s05.id2770682}

\item 共享库（Shared Object，或者Shared Library）\BookAnchor{db.ch18s05.id2770693}
\end{itemize}


共享库留到\hyperlink{db.ch20s04.id2789683}{第 4 节 “共享库”}再详细介绍，本节我们以\hyperlink{db.ch18s03.id2769965}{例 18.\BookSplit{}2 “求一组数的最大值的汇编程序”}为例讨论目标文件和可执行文件的格式。现在详细解释一下这个程序的汇编、链接、运行过程：


\begin{enumerate}

\item 写一个汇编程序保存成文本文件\BookCode{max.\BookSplit{}s}。

\item 汇编器读取这个文本文件转换成目标文件\BookCode{max.\BookSplit{}o}，目标文件由若干个Section组成，我们在汇编程序中声明的\BookCode{.\BookSplit{}section}会成为目标文件中的Section，此外汇编器还会自动添加一些Section（比如符号表）。

\item 然后链接器把目标文件中的Section合并成几个Segment\BookAnchor{db.ch18s05.id2770764}\BookAnchor{db.ch18s05.id2770769}\footnote[28]{\BookAnchor{db.ch18s05.ftn.id2770769}\BookFootnoteBack{db.ch18s05.id2770769}Segment也可以翻译成“段”，为了避免混淆，在本书中只把Section称为段，而Segment直接用英文。}，生成可执行文件\BookCode{max}。

\item 最后加载器（Loader）\BookAnchor{db.ch18s05.id2770791}根据可执行文件中的Segment信息加载运行这个程序。
\end{enumerate}


ELF格式提供了两种不同的视角，链接器把ELF文件看成是Section的集合，而加载器把ELF文件看成是Segment的集合。如下图所示。


\BookFigure{ELF文件}{\BookImage{images/asm.elfoverview.png}{0.529}}{\BookAnchor{db.ch18s05.id2770810}}


左边是从链接器的视角来看ELF文件，开头的ELF Header描述了体系结构和操作系统等基本信息，并指出Section Header Table和Program Header Table在文件中的什么位置，Program Header Table在链接过程中用不到，所以是可有可无的，Section Header Table中保存了所有Section的描述信息，通过Section Header Table可以找到每个Section在文件中的位置。右边是从加载器的视角来看ELF文件，开头是ELF Header，Program Header Table中保存了所有Segment的描述信息，Section Header Table在加载过程中用不到，所以是可有可无的。从上图可以看出，一个Segment由一个或多个Section组成，这些Section加载到内存时具有相同的访问权限。有些Section只对链接器有意义，在运行时用不到，也不需要加载到内存，那么就不属于任何Segment。注意Section Header Table和Program Header Table并不是一定要位于文件的开头和结尾，其位置由ELF Header指出，上图这么画只是为了清晰。



目标文件需要链接器做进一步处理，所以一定有Section Header Table；可执行文件需要加载运行，所以一定有Program Header Table；而共享库既要加载运行，又要在加载时做动态链接，所以既有Section Header Table又有Program Header Table。


\subsection{目标文件}
\BookAnchor{db.ch18s05.id2770854}%


下面用\BookCode{readelf}工具读出目标文件\BookCode{max.\BookSplit{}o}的ELF Header和Section Header Table，然后我们逐段分析。


\BookTerminal{code/ch18s05/01.txt}{9}{22}


ELF Header中描述了操作系统是UNIX，体系结构是80386。Section Header Table中有8个Section Header，从文件地址200（0xc8）开始，每个Section Header占40字节，共320字节，到文件地址0x207结束。这个目标文件没有Program Header。文件地址是这样定义的：文件开头第一个字节的地址是0，然后每个字节占一个地址。


\BookTerminal{code/ch18s05/02.txt}{9}{20}


从Section Header中读出各Section的描述信息，其中\BookCode{.\BookSplit{}text}和\BookCode{.\BookSplit{}data}是我们在汇编程序中声明的Section，而其它Section是汇编器自动添加的。\BookCode{Addr}是这些段加载到内存中的地址（我们讲过程序中的地址都是虚拟地址），加载地址要在链接时填写，现在空缺，所以是全0。\BookCode{Off}和\BookCode{Size}两列指出了各Section的文件地址，比如\BookCode{.\BookSplit{}data}段从文件地址0x60开始，一共0x38个字节，回去翻一下程序，\BookCode{.\BookSplit{}data}段定义了14个4字节的整数，一共是56个字节，也就是0x38。根据以上信息可以描绘出整个目标文件的布局。


\begingroup
\setlength{\tabcolsep}{6pt}
\small
\sloppy
\renewcommand{\arraystretch}{1.16}
\setlength{\parskip}{0pt}
\begin{longtable}{L{\dimexpr 0.386555\linewidth-2\tabcolsep\relax}L{\dimexpr 0.613445\linewidth-2\tabcolsep\relax}}
\BookTableTitle{目标文件的布局}\BookTableLabel{\BookAnchor{db.ch18s05.id2771004}}
\toprule
{\bfseries 起始文件地址} & {\bfseries Section或Header} \\
\midrule
\endfirsthead
\toprule
{\bfseries 起始文件地址} & {\bfseries Section或Header} \\
\midrule
\endhead
0 & ELF Header \\
\addlinespace[2pt]
0x34 & \BookCode{.\BookSplit{}text} \\
\addlinespace[2pt]
0x60 & \BookCode{.\BookSplit{}data} \\
\addlinespace[2pt]
0x98 & \BookCode{.\BookSplit{}bss}（此段为空） \\
\addlinespace[2pt]
0x98 & \BookCode{.\BookSplit{}shstrtab} \\
\addlinespace[2pt]
0xc8 & Section Header Table \\
\addlinespace[2pt]
0x208 & \BookCode{.\BookSplit{}symtab} \\
\addlinespace[2pt]
0x288 & \BookCode{.\BookSplit{}strtab} \\
\addlinespace[2pt]
0x2b0 & \BookCode{.\BookSplit{}rel.\BookSplit{}text} \\
\bottomrule
\end{longtable}
\endgroup


这个文件不大，我们直接用\BookCode{hexdump}工具把目标文件的字节全部打印出来看。


\BookTerminal{code/ch18s05/03.txt}{9}{46}


左边一列是文件地址，中间是每个字节的十六进制表示，右边是把这些字节解释成ASCII码所对应的字符。中间有一个*号表示省略的部分全是0。\BookCode{.\BookSplit{}data}段对应的是这一块：


\BookTerminal{code/ch18s05/04.txt}{8}{6}


\BookCode{.\BookSplit{}data}段将被原封不动地加载到内存中，下一小节会看到\BookCode{.\BookSplit{}data}段被加载到内存地址0x080490a0\textasciitilde{}\BookSplit{}0x080490d7。



\BookCode{.\BookSplit{}shstrtab}和\BookCode{.\BookSplit{}strtab}这两个Section中存放的都是ASCII码：


\BookTerminal{code/ch18s05/05.txt}{9}{10}


可见\BookCode{.\BookSplit{}shstrtab}段保存着各Section的名字，\BookCode{.\BookSplit{}strtab}段保存着程序中用到的符号的名字。每个名字都是以\BookCode{'\textbackslash{}0'}结尾的字符串。



我们知道，C语言的全局变量如果在代码中没有初始化，就会在程序加载时用0初始化。这种数据属于\BookCode{.\BookSplit{}bss}段，在加载时它和\BookCode{.\BookSplit{}data}段一样都是可读可写的数据，但是在ELF文件中\BookCode{.\BookSplit{}data}段需要占用一部分空间保存初始值，而\BookCode{.\BookSplit{}bss}段则不需要。也就是说，\BookCode{.\BookSplit{}bss}段在文件中只占一个Section Header而没有对应的Section，程序加载时\BookCode{.\BookSplit{}bss}段占多大内存空间在Section Header中描述。在我们这个例子中没有用到\BookCode{.\BookSplit{}bss}段，在\hyperlink{db.ch19s03.id2777491}{第 3 节 “变量的存储布局”}会看到这样的例子。



我们继续分析\BookCode{readelf}输出的最后一部分，是从\BookCode{.\BookSplit{}rel.\BookSplit{}text}和\BookCode{.\BookSplit{}symtab}这两个Section中读出的信息。


\BookTerminal{code/ch18s05/06.txt}{9}{20}


\BookCode{.\BookSplit{}rel.\BookSplit{}text}告诉链接器指令中的哪些地方需要做重定位，在下一小节详细讨论。



\BookCode{.\BookSplit{}symtab}是符号表。\BookCode{Ndx}列是每个符号所在的Section编号，例如符号\BookCode{data\_\BookSplit{}items}在第3个Section里（也就是\BookCode{.\BookSplit{}data}段），各Section的编号见Section Header Table。\BookCode{Value}列是每个符号所代表的地址，在目标文件中，符号地址都是相对于该符号所在Section的相对地址，比如\BookCode{data\_\BookSplit{}items}位于\BookCode{.\BookSplit{}data}段的开头，所以地址是0，\BookCode{\_\BookSplit{}start}位于\BookCode{.\BookSplit{}text}段的开头，所以地址也是0，但是\BookCode{start\_\BookSplit{}loop}和\BookCode{loop\_\BookSplit{}exit}相对于\BookCode{.\BookSplit{}text}段的地址就不是0了。从\BookCode{Bind}这一列可以看出\BookCode{\_\BookSplit{}start}这个符号是\BookCode{GLOBAL}的，而其它符号是\BookCode{LOCAL}的，\BookCode{GLOBAL}符号是在汇编程序中用\BookCode{.\BookSplit{}globl}指示声明过的符号。



现在剩下\BookCode{.\BookSplit{}text}段没有分析，\BookCode{objdump}工具可以把程序中的机器指令反汇编（Disassemble）\BookAnchor{db.ch18s05.id2771571}，那么反汇编的结果是否跟原来写的汇编代码一模一样呢？我们对比分析一下。


\BookTerminal{code/ch18s05/07.txt}{9}{25}


左边是机器指令的字节，右边是反汇编结果。显然，所有的符号都被替换成地址了，比如\BookCode{je 23}，注意没有加\BookCode{\$}的数表示内存地址，而不表示立即数。这条指令后面的\BookCode{\textless{}loop\_\BookSplit{}exit\textgreater{}}并不是指令的一部分，而是反汇编器从\BookCode{.\BookSplit{}symtab}和\BookCode{.\BookSplit{}strtab}中查到的符号名称，写在后面是为了有更好的可读性。目前所有指令中用到的符号地址都是相对地址，下一步链接器要修改这些指令，把其中的地址都改成加载时的内存地址，这些指令才能正确执行。


\subsection{可执行文件}
\BookAnchor{db.ch18s05.id2771639}%


现在我们按上一节的步骤分析可执行文件\BookCode{max}，看看链接器都做了什么改动。


\BookTerminal{code/ch18s05/08.txt}{9}{67}


在ELF Header中，\BookCode{Type}改成了\BookCode{EXEC}，由目标文件变成可执行文件了，\BookCode{Entry point ad\BookSplit{}dress}改成了0x8048074（这是\BookCode{\_\BookSplit{}start}符号的地址），还可以看出，多了两个Program Header，少了两个Section Header。



在Section Header Table中，\BookCode{.\BookSplit{}text}和\BookCode{.\BookSplit{}data}段的加载地址分别改成了0x08048074和0x080490a0。\BookCode{.\BookSplit{}bss}段没有用到，所以被删掉了。\BookCode{.\BookSplit{}rel.\BookSplit{}text}段就是用于链接过程的，做完链接就没用了，所以也删掉了。



多出来的Program Header Table描述了两个Segment的信息。\BookCode{.\BookSplit{}text}段和前面的ELF Header、Program Header Table一起组成一个Segment（\BookCode{FileSiz}指出总长度是0x9e），\BookCode{.\BookSplit{}data}段组成另一个Segment（总长度是0x38）。\BookCode{VirtAddr}列指出第一个Segment加载到虚拟地址0x08048000（注意在x86平台上后面的\BookCode{PhysAddr}列是没有意义的，并不代表实际的物理地址），第二个Segment加载到地址0x080490a0。\BookCode{Flg}列指出第一个Segment的访问权限是可读可执行，第二个Segment的访问权限是可读可写。最后一列\BookCode{Align}的值0x1000（4K）是x86平台的内存页面大小。在加载时文件也要按内存页面大小分成若干页，文件中的一页对应内存中的一页，对应关系如下图所示。


\BookFigure{文件和加载地址的对应关系}{\BookImage{images/asm.load.png}{0.668}}{\BookAnchor{db.ch18s05.id2771830}}


这个可执行文件很小，总共也不超过一页大小，但是两个Segment必须加载到内存中两个不同的页面，因为MMU的权限保护机制是以页为单位的，一个页面只能设置一种权限。此外还规定每个Segment在文件页面内偏移多少加载到内存页面仍然要偏移多少，比如第二个Segment在文件中的偏移是0xa0，在内存页面0x08049000中的偏移仍然是0xa0，所以从0x080490a0开始，这样规定是为了简化链接器和加载器的实现。从上图也可以看出\BookCode{.\BookSplit{}text}段的加载地址应该是\BookCode{0x08048074}，\BookCode{\_\BookSplit{}start}符号位于\BookCode{.\BookSplit{}text}段的开头，所以\BookCode{\_\BookSplit{}start}符号的地址也是0x08048074，从符号表中可以验证这一点。



原来目标文件符号表中的\BookCode{Value}都是相对地址，现在都改成绝对地址了。此外还多了三个符号\BookCode{\_\BookSplit{}\_\BookSplit{}bss\_\BookSplit{}start}、\BookCode{\_\BookSplit{}edata}和\BookCode{\_\BookSplit{}end}，这些符号在链接脚本中定义，被链接器添加到可执行文件中，链接脚本在\hyperlink{db.ch20s01.id2786161}{第 1 节 “多目标文件的链接”}介绍。



再看一下反汇编的结果：


\BookTerminal{code/ch18s05/09.txt}{9}{25}


指令中的相对地址都改成绝对地址了。我们仔细检查一下改了哪些地方。首先看跳转指令，原来目标文件的指令是这样：


\BookTerminal{code/ch18s05/10.txt}{9}{7}


现在改成了这样：


\BookTerminal{code/ch18s05/11.txt}{9}{7}


改了吗？其实只是反汇编的结果不同了，指令的机器码根本没变。为什么不用改指令就能跳转到新的地址呢？因为跳转指令中指定的是相对于当前指令向前或向后跳多少字节，而不是指定一个完整的内存地址，内存地址有32位，这些跳转指令只有16位，显然也不可能指定一个完整的内存地址，这称为相对跳转。这种相对跳转指令只有16位，只能在当前指令前后的一个小范围内跳转，不可能跳得太远，也有的跳转指令指定一个完整的内存地址，可以跳到任何地方，这称绝对跳转，在\hyperlink{db.ch20s04.id2790795}{第 4.\BookSplit{}2 节 “动态链接的过程”}我们会看到这样的例子。



再看内存访问指令，原来目标文件的指令是这样：


\BookTerminal{code/ch18s05/12.txt}{7}{5}


现在改成了这样：


\BookTerminal{code/ch18s05/13.txt}{7}{5}


指令中的地址原本是0x00000000，现在改成了0x080409a0（注意是小端字节序）。那么链接器怎么知道要改这两处呢？是根据目标文件中的\BookCode{.\BookSplit{}rel.\BookSplit{}text}段提供的重定位信息来改的：


\BookTerminal{code/ch18s05/14.txt}{8}{6}


第一列\BookCode{Offset}的值就是\BookCode{.\BookSplit{}text}段需要改的地方，在\BookCode{.\BookSplit{}text}段中的相对地址是8和0x17，正是这两条指令中00 00 00 00的位置。

\BookNav{\hyperlink{db.ch18s04}{上一页}\enspace\textperiodcentered\enspace \hyperlink{db.ch18}{上一级}\enspace\textperiodcentered\enspace \hyperlink{db.index}{起始页}\enspace\textperiodcentered\enspace \hyperlink{db.ch19}{下一页}}
